%% C4 nivel 1 — Diagrama de contexto (plantilla)
%% Compilar: tectonic -X compile c4-contexto.tex
\documentclass[11pt,a4paper]{article}
\usepackage{../../docs/latex/legasoft}
\usepackage{../../docs/latex/legasoft-diagramas}

\lgtitulo{C4 nivel 1: diagrama de contexto}
\lgtituloCorto{C4 — Contexto}
\lgcodigo{LGS-MOD-C4-1}
\lgversion{0.1}
\lgestado{Plantilla}
\lgfecha{\campo{fecha}}
\lgautor{\lgArquitecto\ (\lgArquitectoRol)}

\begin{document}

\begin{center}
{\sffamily\bfseries\LARGE\color{lgPrimario} C4 nivel 1: diagrama de contexto\par}
\vspace{2pt}
{\sffamily\normalsize\color{lgGris} \campo{nombre del sistema} · Legasoft\par}
\end{center}
\vspace{6pt}
{\color{lgBorde}\hrule height 0.8pt}
\vspace{10pt}

\begin{porcompletar}[Cómo usar esta plantilla]
El nivel 1 muestra el sistema como una \emph{caja negra} y su relación con
personas y sistemas externos. No incluye tecnología interna ni contenedores: eso
corresponde al nivel 2. Sustituya los rótulos en ámbar por los elementos reales
del proyecto y elimine esta caja.
\end{porcompletar}

% =============================================================================
\section*{Diagrama}

\begin{figure}[H]
\centering
\ajustaancho{%
\begin{tikzpicture}[c4]

  % --- Personas -------------------------------------------------------------
  \node[c4persona] (usuario)
    {\ctexto{\campoclaro{Usuario operativo}}{}{Registra y consulta información en el día a día.}};

  \node[c4persona, right=3.4cm of usuario] (gestor)
    {\ctexto{\campoclaro{Responsable de gestión}}{}{Consulta reportes y toma decisiones.}};

  % --- Sistema en alcance ---------------------------------------------------
  \node[c4focoancho, below=2.3cm of $(usuario)!0.5!(gestor)$]
    (sistema)
    {\ctexto{\campo{Nombre del sistema}}{}{Digitaliza y controla
     \campo{el proceso objetivo}.}};

  % --- Sistemas externos ----------------------------------------------------
  \node[c4externo, below left=2.4cm and 0.4cm of sistema.south] (ext1)
    {\ctexto{\campo{Sistema externo A}}{}{\campo{Qué provee o consume.}}};

  \node[c4externo, below right=2.4cm and 0.4cm of sistema.south] (ext2)
    {\ctexto{\campo{Sistema externo B}}{}{\campo{Qué provee o consume.}}};

  % --- Relaciones -----------------------------------------------------------
  \draw[c4rel] (usuario.south) -- (usuario.south |- sistema.north)
    node[c4etq, pos=0.5] {\campo{registra datos}\\\campo{HTTPS}};

  \draw[c4rel] (gestor.south) -- (gestor.south |- sistema.north)
    node[c4etq, pos=0.5] {\campo{consulta reportes}\\\campo{HTTPS}};

  \draw[c4relbi] (sistema.south) -- (ext1.north -| ext1.north)
    node[c4etq, pos=0.55] {\campo{intercambia}\\\campo{API REST}};

  \draw[c4rel] (sistema.south) -- (ext2.north)
    node[c4etq, pos=0.55] {\campo{envía}\\\campo{protocolo}};

\end{tikzpicture}}
\caption{Contexto del sistema: actores, sistema en alcance y sistemas externos.}
\end{figure}

\cleyenda

% =============================================================================
\Needspace*{20\baselineskip}
\section*{Elementos del diagrama}

\begin{center}
\begin{tabularx}{\linewidth}{@{}L{3.8cm} L{2.6cm} Y@{}}
\toprule
\sffamily\bfseries\color{lgPrimario}Elemento &
\sffamily\bfseries\color{lgPrimario}Tipo &
\sffamily\bfseries\color{lgPrimario}Descripción y responsabilidad \\
\midrule
\campo{Usuario operativo}     & Persona          & \campo{qué hace con el sistema.} \\
\campo{Responsable de gestión}& Persona          & \campo{qué necesita del sistema.} \\
\campo{Nombre del sistema}    & Sistema en alcance & \campo{propósito principal.} \\
\campo{Sistema externo A}     & Sistema externo  & \campo{qué información intercambia y por qué.} \\
\campo{Sistema externo B}     & Sistema externo  & \campo{qué información intercambia y por qué.} \\
\bottomrule
\end{tabularx}
\captionof{table}{Inventario de elementos del diagrama de contexto.}
\end{center}

\begin{nota}[Siguiente nivel]
El detalle interno del sistema se describe en \ruta{c4-contenedores.tex}
(nivel 2) y la estructura de cada contenedor en \ruta{c4-componentes.tex}
(nivel 3). La descripción narrativa acompaña al documento de arquitectura
\texttt{LGS-ARQ-001}.
\end{nota}

\end{document}
